% The style file loads xcolor itself, so the table option has to be handed over
% before that happens. It brings in colortbl, for row and cell backgrounds.
% Table 4 used it for a reversed-out legend band until 2026-08-23; nothing in
% the paper draws on a colour now, so this can go if nothing new wants it.
\PassOptionsToPackage{table}{xcolor}
\usepackage{style/usenix-2020-09}

\usepackage{tikz}
\usepackage{helvet}
% arrows.meta is loaded alongside arrows, not instead of it: the older figures
% keep their >=stealth' tips, while fig-thread-race sizes its heads in absolute
% pt. Old style tips scale with the line width, which made the 3pt connectors
% in that figure read as blobs rather than arrows.
\usetikzlibrary{arrows,arrows.meta,backgrounds,decorations.pathreplacing,patterns.meta}
\usepackage{amsmath}
\usepackage{amssymb}
\usepackage{bm}
\usepackage{graphicx}
\usepackage{booktabs}
\usepackage{array}
\usepackage{algorithm}
\usepackage{algpseudocode}
\usepackage{float}
\usepackage{makecell}
\usepackage{comment}

% The bibliography carries long URLs and \texttt{} commit hashes. In a narrow
% two-column measure those are unbreakable boxes, so TeX pads the inter-word
% spaces around them instead, which shows up as rivers of white in an entry.
% xurl lets a URL break at any character; emergencystretch gives the paragraph
% builder a last-resort pass so it stops stretching good lines to save a bad one.
\usepackage{xurl}

% breaklinks, because the usenix style loads hyperref with no options and the
% pdftex default is false. A link that breaks across lines then gets ONE
% bounding box spanning everything between its start and its end, and xurl
% breaks long URLs freely, so a citation whose URL straddled a column break
% produced a clickable rectangle 200pt tall covering the whole of Figure 5.
% With breaklinks the link is emitted one rectangle per line.
\hypersetup{breaklinks=true}
